% Evidence: canonical Phase 5B condition_metrics.json, paired_contrasts.json,
% uncertainty_results.json, and sensitivity_analysis.csv.
\begin{table*}[t]
\centering
\begin{tabular}{@{}lrrrrr@{}}
\toprule
Condition & Mean regret & Invalid & Harmful & Claim valid & Failed \\
\midrule
A & 4.723122 & 38.89\% & 16.67\% & 61.11\% & 1 \\
B & 1.389345 & 11.11\% & 11.11\% & 88.89\% & 1 \\
C & 0.250287 & 0\% & 0\% & 100\% & 0 \\
D & 3.598322 & 22.22\% & 11.11\% & 77.78\% & 0 \\
\midrule
\multicolumn{6}{@{}l}{\textit{Paired all-cell regret contrasts
(scenario-cluster interval)}}\\
\(B-A\) & \multicolumn{2}{l}{\(-3.333777\)}
  & \multicolumn{3}{l}{\([-4.111998,-2.555556]\)} \\
\(C-B\) & \multicolumn{2}{l}{\(-1.139058\)}
  & \multicolumn{3}{l}{\([-2.444444,0.166328]\)} \\
\(D-C\) & \multicolumn{2}{l}{\(+3.348035\)}
  & \multicolumn{3}{l}{\([1.807181,4.888889]\)} \\
\bottomrule
\end{tabular}
\caption{Canonical fail-closed outcomes, \(N=18\) per condition.  Lower
regret is better.  Intervals use only two scenario clusters and are
pilot-descriptive, not conventional significance claims.}
\label{tab:main-results}
\end{table*}
